%!TeX program = xelatex
\documentclass[12pt,a4paper,UTF8]{ctexart}
\usepackage{gbt7714}
\usepackage{zjsureport}

\begin{document}

\cover
\maketitle
\begin{abstract}
此处填写摘要正文。摘要环境现在支持普通段落自动换行，直接连续输入中文内容即可，不需要手动插入换行命令；如果摘要较长，LaTeX 会按照当前版心宽度自动断行并分页。
\end{abstract}
\keywords{机器学习；精准营销；XGBoost}

\clearpage
\tableofcontents

\clearpage
\section{引言}

\subsection{研究背景}
此处填写研究背景。\cite{gu2020xgboost}

\subsection{研究意义}
此处填写研究意义。这里的正文已经恢复为正常 UTF-8 文本，修改标题后重新编译，目录会与正文同步更新。

\section{数据准备}

\subsection{数据描述}
此处填写数据描述。

\subsection{数据预处理}

\subsubsection{剔除无关变量}
此处填写剔除无关变量的说明。

\subsubsection{缺失值检查}
此处填写缺失值检查的说明。

\subsubsection{重复值检查}
此处填写重复值检查的说明。

\subsubsection{异常值检查}
此处填写异常值检查的说明。

\subsection{描述性统计}

\subsubsection{连续变量数据情况}
此处填写连续变量统计说明。

\subsubsection{分类变量数据情况}
此处填写分类变量统计说明。

\section{特征提取}

\subsection{One-Hot编码}
此处填写 One-Hot 编码说明。

\subsection{数据归一化}
此处填写数据归一化说明。

\section{模型的构建}

\subsection{模型选择}
此处填写模型选择说明。

\subsection{模型评价}
此处填写模型评价说明。

\section{总结}
此处填写总结内容。

\reference

\clearpage
\section*{附录}
\addcontentsline{toc}{section}{附录}
\begin{lstlisting}
#####导入数据#####
value <- read.table(file = file.choose(), header = TRUE, sep = ",")
data <- value

#####数据预处理#####
#####1、删除没用的列变量#####
data1 <- data[, -c(14:17)]

#####2、查看数据中是否有缺失值#####
list <- list()
length(list) = 14
for (i in 1:14) { list[[i]] = table(is.na(data1[, i])) }
unlist(list)

#####3、查看数据中是否有重复值#####
table(duplicated(data1$ID))

#####4、查看数据中是否存在异常值#####
par(mfcol = c(2, 3))
par(mai = c(0, 0.5, 0.5, 0))
par(oma = c(3, 0, 2, 2))
boxplot(data1$age)
mtext("age", side = 1, line = 1, outer = FALSE)
boxplot(data1$duration)
mtext("duration", side = 1, line = 1, outer = FALSE)
boxplot(data1$campaign)
mtext("campaign", side = 1, line = 1, outer = FALSE)
boxplot(data1$previous)
mtext("previous", side = 1, line = 1, outer = FALSE)
boxplot(data1$balance)
mtext("balance", side = 1, line = 1, outer = FALSE)
mtext("异常值检查", side = 3, line = -2, outer = TRUE, cex = 1.5)

#####描述性统计#####
#####1、连续变量描述性统计#####
list <- list()
length(list) = 5
lianxu <- cbind(data1$age, data1$duration, data1$campaign, data1$previous, data1$balance)
for (i in 1:5) {
  list[[i]] = fivenum(lianxu[, i])
}
mean <- numeric(5)
for (i in 1:5) { mean[i] = mean(lianxu[, i]) }

#####2、分类变量描述性统计#####
par(mfcol = c(3, 2))
par(mai = c(0.5, 0.5, 0.2, 0))
par(oma = c(2, 0, 3, 2))
tu1 <- table(as.vector(data1$marital)[c(which(data1$y == "n"))])
tu2 <- table(as.vector(data1$marital)[c(which(data1$y == "y"))])
tu <- cbind(as.data.frame(tu1), Freq2 = as.data.frame(tu2)[, 2])[, 2:3]
rownames(tu) <- c(as.character(as.data.frame(tu1)[, 1]))
barplot(as.matrix(t(tu)), beside = TRUE)
\end{lstlisting}

\end{document}
